\section{Escalado en tiempo de prueba: de los LLM a los modelos difusos}

\subsection{El fin de la era del preentrenamiento y el cómputo en inferencia}

Ilya Sutskever (cofundador de OpenAI) planteó una idea influyente: la era del preentrenamiento está llegando a su fin. Para los grandes modelos de lenguaje (LLM), los datos de entrenamiento disponibles en internet se están agotando. Incluso si se sigue aumentando la escala del modelo o la cantidad de cómputo de entrenamiento, las mejoras en el rendimiento serán cada vez más limitadas.

\begin{knowledgebox}{El surgimiento del escalado en tiempo de prueba}
La idea central del escalado en tiempo de prueba (test-time scaling) es: en lugar de invertir más recursos computacionales durante el entrenamiento, es mejor gastar más cómputo durante la inferencia para mejorar la calidad de la salida. Esta idea ya ha obtenido un éxito significativo en el campo de los LLM:
\begin{itemize}
    \item Los modelos O1/O3 de OpenAI invierten más tiempo durante la inferencia mediante un "modo de pensamiento" para mejorar las respuestas.
    \item DeepSeek R1 entrena con RL (aprendizaje por refuerzo) al modelo para que genere tokens de pensamiento intermedios, logrando un "momento aha" —el modelo puede descubrir y corregir errores de razonamiento por sí mismo—.
    \item En el benchmark ARC-AGI, la puntuación de O3 saltó del 10\%--20\% de los modelos anteriores hasta alcanzar el 88\%.
\end{itemize}
\end{knowledgebox}

\subsection{La transición desde LLM hasta modelos difusos}

La pregunta central del curso es: ¿es posible trasladar la idea del escalado en tiempo de prueba desde los LLM hacia los modelos difusos y los modelos de flujo? El ponente considera que este es un enfoque muy prometedor. Además, las técnicas en tiempo de inferencia requieren menos recursos computacionales que el cómputo en entrenamiento, lo que permite a equipos de investigación pequeños desarrollar modelos de alto rendimiento.

\begin{importantbox}{El marco Proposer-Verifier}
La arquitectura básica del escalado en tiempo de prueba se compone de dos componentes:
\begin{itemize}
    \item \textbf{Proposer (proponente)}: un modelo generativo preentrenado (como un modelo difuso), encargado de generar salidas candidatas.
    \item \textbf{Verifier (verificador)}: otro modelo o función que evalúa la calidad de la salida del proposer. Puede ser una red neuronal, un verificador de código o cualquier función de puntuación.
\end{itemize}
Basándose en el feedback del verifier, se puede (1) guiar el proceso de generación del proposer o (2) hacer que el proposer genere múltiples candidatos para seleccionar el óptimo. Esto es análogo al marco "generador-discriminador" de las GANs.
\end{importantbox}

\subsection{Guía en tiempo de prueba desde la perspectiva del RL}

La presentación de Dan (ingeniero de OpenAI) ilustra con viveza la visión del escalado en tiempo de prueba: al igual que a Einstein le tomaron 8 años desarrollar la relatividad general en 1907, el modelo O3 solo necesita un minuto de "pensamiento" para resolver un examen final de relatividad general —lo cual depende de invertir más cómputo durante la inferencia—.

Yann LeCun alguna vez comparó el preentrenamiento con una "torta grande", y el aprendizaje por refuerzo (RL) era simplemente las "cerezas de la torta". Sin embargo, la hoja de ruta de OpenAI indica que en el futuro el cómputo de RL ocupará una proporción cada vez mayor del cómputo total, superando potencialmente al cómputo de preentrenamiento.

\begin{warningbox}{Implicaciones de las tendencias de investigación}
El ponente observa que la comunidad actual de investigación en IA ya no se centra únicamente en desarrollar nuevas tecnologías individuales o publicar artículos. Cada vez más equipos de investigación (incluidas startups) están aprovechando las técnicas de escalado en tiempo de prueba para superar el rendimiento de los modelos existentes utilizando recursos computacionales relativamente menores en comparación con entrenar un modelo grande desde cero. Esto ofrece una oportunidad importante a laboratorios académicos y pequeños equipos con recursos limitados: no es necesario contar con miles de GPUs; también se pueden obtener resultados competitivos mediante técnicas ingeniosas en tiempo de inferencia.
\end{warningbox}

\subsection{Resumen de la sección}

El escalado en tiempo de prueba ya ha demostrado un gran potencial en el campo de los LLM. Trasladar esta idea hacia los modelos difusos, mejorando la calidad de la generación durante la inferencia a través del marco proposer-verifier, es una dirección importante en la investigación actual.


